\documentclass[12pt, a4paper]{article}

\usepackage[utf8]{inputenc}   % für deutsche Umlaute (input)
\usepackage[ngerman]{babel}   % für neue deutsche Rechtschreibung 
\usepackage[T1]{fontenc}      % für deutsche Umlaute (output)

\usepackage{amsmath}          % für mathematische Formeln
\usepackage{amssymb}          % für mathematische Symbole

\usepackage{graphicx}         % für Bilder

\usepackage[numbers]{natbib}
\bibliographystyle{plainnat}



\setlength{\parindent}{0em}   % Einzug (Einrückung) der ersten Zeile im Abschnitt
\setlength{\parskip}{0.75em}  % Abstand zwischen den einzelnen Abschnitten

\title{\LaTeX-Intro}          % Titel der Arbeit, wird für \maketitle benötigt
\author{Peter Kessler}        % Autor der Arbeit, wird für \maketitle benötigt
\date{März 2019}              % Datum der Arbeit, wird für \maketitle benötigt


\begin{document}

\maketitle

\pagebreak
\tableofcontents


\pagebreak 
\section{Herkunft}
Das Basis-Programm von LaTeX ist TeX und wurde von Donald E. Knuth während seiner Zeit als Informatik-Professor an der Stanford University entwickelt. Auf TeX aufbauend entwickelte Leslie Lamport\cite{Lamport2017} Anfang der 1980er Jahre LaTeX, eine Sammlung von TeX-Makros, die die Benutzung für den durchschnittlichen Anwender gegenüber TeX vereinfachten und erweiterten. Der Name LaTeX ist eine Abkürzung für Lamport TeX.

\begin{figure}[h!]
\centering
  \includegraphics[width=0.5\textwidth]{./Bilder/overleaf.png}
  \caption{Overleaf - \LaTeX in der Cloud}
\end{figure}

Ein sehr zu empfehlendes Buch zum Thema \LaTeX\ ist\cite{Kottwitz2011}. 



\pagebreak
\section{WYSIWYG vs WYSIWYAF}

\subsection{WYSIWYG}
Im Gegensatz zu anderen Textverarbeitungsprogrammen, die nach dem WYSIWYG (What-you-see-is-what-you-get) Prinzip funktionieren, arbeitet der Autor mit Textdateien, in denen er innerhalb eines Textes anders zu formatierende Passagen oder Überschriften mit Befehlen textuell auszeichnet.

\subsection{WYSIWYAF}
Das dabei von LaTeX generierte Layout gilt als sehr sauber, sein Formelsatz als sehr ausgereift. Außerdem ist die Ausgabe u. a. nach PDF, HTML und PostScript möglich. LaTeX eignet sich insbesondere für umfangreiche Arbeiten wie Diplomarbeiten und Dissertationen, die oftmals strengen typographischen Ansprüchen genügen müssen. Insbesondere in der Mathematik und den Naturwissenschaften erleichtert LaTeX das Anfertigen von Schriftstücken durch seine komfortablen Möglichkeiten der Formelsetzung gegenüber üblichen Textverarbeitungssystemen. Das Verfahren von LaTeX wird auch mit WYSIWYAF (What you see is what you asked for.) umschrieben.

Wie TeX selbst ist LaTeX weitestgehend rechnerunabhängig verwendbar. Das bedeutet, dass es für die meisten Betriebssysteme analog zu TeX auch für LaTeX lauffähige, produktiv einsetzbare LaTeX-Installationen gibt. Zu diesen Betriebssystemen gehören zum Beispiel Microsoft Windows von der Version 3.x bis zur aktuellen Version 10, Apple macOS sowie diverse Linux-Distributionen. Unter der Voraussetzung, dass alle verwendeten Zusatzpakete (siehe unten) in geeigneten Versionen installiert sind, besteht der Vorteil der Verwendung von LaTeX darin, dass das Ergebnis unabhängig von der verwendeten Rechnerplattform und dem verwendeten Drucker in den beiden Ausgabeformaten DVI und PDF hinsichtlich der Schriftarten und -größen sowie der Zeilen- und Seitenumbrüche im Druck stets exakt gleich ist. LaTeX ist hierbei nicht auf die Schriftarten des jeweiligen Betriebssystems angewiesen. Jene Betriebssystem-Schriftarten sind häufig für die Anzeige am Monitor optimiert. LaTeX enthält eine Reihe eigener Schriftarten, die ihrerseits für den Druck optimiert sind. 




\pagebreak
\section{Arbeiten mit mathematischen Formeln}
LaTeX unterstützt wie kaum ein anderes Textsatzungsprogramm das Einfügen mathematischer Formeln. Wahlweise können Formeln innerhalb einer Absatzes („Inline“) oder als eigenständiger Absatz in ein Dokument eingefügt werden.



\subsection{Inline Formeln}
Pythagoras sagt: Seien $a$ und $b$ die 
Katheten und $c$ die Hypotenuse, dann gilt: $a^2+b^2=c^2$. 
Somit gilt für die Hypothenuse: $c=\sqrt{a^2+b^2}$.



\subsection{Abgesetzte Formeln}

Pythagoras sagt: Seien $a$ und $b$ die Katheten und $c$ die
Hypotenuse, dann gilt: 

\begin{displaymath}
  a^2+b^2=c^2 
\end{displaymath}
 
Somit gilt für die Hypothenuse: 

\begin{displaymath}
  c=\sqrt{a^2+b^2} 
\end{displaymath}


\subsection{Ausgerichtete Formeln}

Da gibt es zwei Formeln, die jedes Kind kennt: 
\begin{align}
  a^2 + b^2 &= c^2     \label{pythagoras}  \\  
  e &= m c^2           \label{einstein}  
\end{align}

Wobei \eqref{pythagoras} Pythagoras und \eqref{einstein} Albert Einstein zugeschrieben wird.



\pagebreak
\section{Zitieren und Literaturverzeichnis}

Quellenangaben und Zitate erfüllen in wissenschaftlichen Texten bzw. in wissenschaftlichen Arbeiten zwei Hauptziele\cite{NW2017}: 

\begin{itemize}

\item[a)] Nachvollziehbarkeit: Ein wissenschaftlicher Beitrag kann nur dann weiterverwendet werden, wenn sich die Argumentation für die Lesenden überprüfen lässt.

\item[b)] Unterscheidung zwischen eigenen und fremden Gedanken/Ideen: Bilder, Grafiken, Tabellen, Texte die aus anderen Arbeiten stammen, müssen als fremdes geistiges Eigentum ausgewiesen werden.

\end{itemize}

\textbf{Wichtige Regel:} Quellenangaben gehören zum Satz dazu und stehen daher VOR dem Punkt bzw. Satzzeichen.



\pagebreak
\addcontentsline{toc}{section}{Literatur-/Quellenverzeichnis}
\renewcommand{\refname}{Literatur-/Quellenverzeichnis} 
\bibliography{./Verzeichnisse/Literatur}

\end{document}
